\chapter{奶蛙的语言—符号系统：内容、形式与立场}\label{ch:06}

\begin{nlkeypoints}
  \begin{enumerate}[label=(\arabic*),leftmargin=2.4em]
    \item 以 Shifman 三维框架解剖奶蛙语料，得到一个\keyw{“一软两硬”}结构：
          内容维度高度可替换（任何“黄色大笑生物”都能接入），
          形式维度高度稳定（结尾集体大笑、衍生语搭配、
          “压力桥段$\to$大笑”的叙事套路），立场维度高度约定化
          （使用奶蛙即声明“我绷不住了，但我选择用笑来处理”）。
    \item 形象谱系遵循\keyw{窄通道规律}：轮廓趋同（黄色主体、大笑姿态、
          膨大腹部），细节发散（换皮、物化、娘化）。
          谱系表按“保留特征 / 丢失特征”两列组织，
          判定依据是第五章迷因簇归并标准 (C1)--(C3)。
    \item “捧腹大笑”是一个\keyw{声音符号}而非单纯的情绪表达：
          它的分析不依赖音源考据是否成立；
          衍生语是\keyw{空耳}（\en{soramimi}）驱动的语音梗，
          其“语音外壳保留、语义外壳被重新指派”的双层结构
          是奶蛙语言系统的核心装置。
  \end{enumerate}
\end{nlkeypoints}

第五章配好了工具，本章开始解剖。解剖对象是问题 Q1：
奶蛙究竟是一个怎样的符号系统？本章的材料是三类语料：
形象语料（谱系与变种）、语言语料（衍生语与固定音轨）、
动作语料（衍生动作与叙事套路）。除正典设定来自 A/B 级来源外，
这三类语料几乎全部来自社群考据与第一版教材的记录（\conf{C} 级），
因此本章的每一张谱系表都附带“判定依据”列，
读者可以按第五章的标准 (C1)--(C3) 自行复核归并是否成立。
本章的结论将是一个结构性判断：奶蛙的符号系统呈
“一软两硬”的不对称——这个判断本身是本书的分析（\conf{M} 级），
但它解释了第八章将要观察到的现象：
为什么 AIGC 时代的二创在形象上千姿百态，
在形式上却高度重复。

\section{分析框架：三维的操作化}\label{sec:06:framework}

Shifman 把数字迷因的分析从“内容像什么”推进到“它在做什么”：
每个迷因用例同时携带\keyw{内容}（\en{content}）、
\keyw{形式}（\en{form}）与\keyw{立场}（\en{stance}）
三个维度\cite{shifman2013memes}。框架本身是 A 级理论资源；
把它落到奶蛙语料上，需要先做一步操作化：

\begin{table}[htbp]\centering
  \caption{三维框架在奶蛙语料上的操作化（判定问题均为可观察问题）}
  \label{tab:06:ops}
  \small
  \begin{tabularx}{\textwidth}{@{}llX@{}}
    \toprule
    维度 & 操作化提问 & 奶蛙语料中的对应物 \\ \midrule
    内容 & “它画的是什么？” &
      形象（黄、膨腹、细肢、黑手）、场景（校园、职场、
      游戏跨界）、被嫁接的文本（课程、考试、品牌） \\
    形式 & “它按什么套路组织？” &
      结尾集体大笑、变调笑声音轨、衍生语搭配、
      “压力桥段$\to$大笑收尾”的叙事结构、竖屏短时长 \\
    立场 & “它在替谁说什么话？” &
      “绷不住但用笑处理”的自陈、自嘲、
      对精致人设的疏离、圈内身份的默示声明 \\ \bottomrule
  \end{tabularx}
\end{table}

操作化的纪律是：每一维的判定问题必须能落到
“看一段二创就能回答”的层面。“立场”最容易被写成印象式的
“反映了当代青年的心态”，本书拒绝这种写法——
立场维度只讨论\keyw{可指认的语用行为}：
用梗图代替自述（第七章机制 M2）、
在何种情境下使用哪个衍生语（本章 \ref{sec:06:form} 节）、
以及社群自己写下的立场声明文本（\ref{sec:06:stance} 节）。

\section{内容维度：形象谱系学}\label{sec:06:content}

\subsection{谱系总表}

奶蛙的形象不是一张图，而是一棵树。按第一版教材的整理与
本书的复核（\conf{C} 级，来源为第一版教材与社群考据词条
\cite{lostmedia2026naiwa}），谱系按“保留多少正典特征”
分为三级：换皮、异形、客串。

\begin{table}[htbp]\centering
  \caption{奶蛙形象谱系（据第一版教材整理，\conf{C} 级；
          “保留 / 丢失”两列即归并标准 (C1) 的操作化）}
  \label{tab:06:pedigree}
  \small
  \begin{tabularx}{\textwidth}{@{}llXX@{}}
    \toprule
    类别 & 代表 & 保留特征 & 丢失 / 替换特征 \\ \midrule
    换皮奶蛙 & 异色奶蛙 & 全部形体结构 &
      仅配色变化 \\
    & “植奶”“僵奶”（《植奶大战僵龙》系列） &
      形体结构、动作程式 &
      皮肤替换为《植物大战僵尸》元素 \\
    & 线代课本奶蛙、VIVO 奶蛙、学习通奶蛙 &
      轮廓与姿态 & 表面替换为图标、封面等平面元素 \\
    异形奶蛙 & 奶蛋 & 眼睛与大笑 &
      躯干、四肢全部简化 \\
    & 奶狗、奶鸡、奶鼠、奶茶、豆奶 &
      眼睛与大笑（部分含黄色） & 物种形态整体替换 \\
    & 奶蛙娘 & 面部特征、黑色双手 & 人体化（娘化） \\
    & 物化奶龙（如“奶菜”） & 局部特征 & 拟物化 \\
    客串角色 & 美团袋鼠 & 姿态与笑声 & 伪装为其他 IP / 品牌形象 \\
    & 牛来、牛来妈 & 姿态与笑声 & 独立角色化 \\ \bottomrule
  \end{tabularx}
\end{table}

\begin{definition}[换皮与异形（本书用法）]\label{def:06:reskin}
  \textbf{换皮奶蛙}：仅替换表面元素（配色、图标、封面），
  保留奶蛙全部形体结构的变种。
  \textbf{异形奶蛙}：丢失大量形体特征，
  仅保留“眼睛与大笑”这一最小辨识单元的变种。
  两者以“丢失特征的比例”划界，
  而非以“像不像青蛙”划界——奶蛙娘虽是人形，
  但因完整保留面部特征与黑色双手，
  在社群辨识实践中归入异形而非独立形象。
\end{definition}

这条分界线来自社群的使用实践（归并标准 (C3)），
不是本书的发明：一个形象只要还能被社群成员
在提示下识别为“这是奶蛙”，它就还在谱系内；
识别一旦失败，谱系就到了边界。
谱系因此是\keyw{辨识边界的图谱}，不是审美相似度的图谱。

\begin{figure}[htbp]\centering
\begin{tikzpicture}[font=\small,
  every node/.style={align=center,inner sep=3pt},
  lv/.style={rounded corners,draw=NLDeep,thick,fill=NLYellow!18,
    minimum height=1.1cm,minimum width=2.6cm},
  leaf/.style={rounded corners=1pt,draw=NLGrey,fill=NLBg,
    minimum height=0.8cm,font=\footnotesize},
  arr/.style={-{Stealth[length=2.6mm]},thick,NLDeep}]
  \node[lv] (canon) {正典奶龙\\{\footnotesize 黄色、圆头圆脑（\conf{A/B}）}};
  \node[lv,below=1.0cm of canon] (naiwa) {经典奶蛙\\{\footnotesize AI 渲染畸变：膨腹、细肢（\conf{C}）}};
  \node[lv,below left=1.1cm and 1.4cm of naiwa] (reskin) {换皮奶蛙};
  \node[lv,below=1.1cm of naiwa] (hetero) {异形奶蛙};
  \node[lv,below right=1.1cm and 1.4cm of naiwa] (guest) {客串角色};
  \node[leaf,below=0.55cm of reskin] (l1) {异色 / 植奶僵奶 / 图标换皮};
  \node[leaf,below=0.55cm of hetero] (l2) {奶蛋 / 奶狗等 / 奶蛙娘 / 物化};
  \node[leaf,below=0.55cm of guest] (l3) {美团袋鼠 / 牛来};
  \draw[arr] (canon) -- node[right,font=\footnotesize]{一次大畸变} (naiwa);
  \draw[arr] (naiwa) -- node[above left,font=\footnotesize]{持续小变异} (reskin);
  \draw[arr] (naiwa) -- (hetero);
  \draw[arr] (naiwa) -- (guest);
  \draw[arr,dashed,NLGrey] (reskin) -- (l1);
  \draw[arr,dashed,NLGrey] (hetero) -- (l2);
  \draw[arr,dashed,NLGrey] (guest) -- (l3);
\end{tikzpicture}
  \caption{奶蛙形象谱系树（据第一版教材整理，\conf{C} 级）。
          实线箭头为“变异方向”，均为自上而下；
          第五章“方向性不对称”在谱系上表现为：
          没有从子类回到父类的箭头。}
  \label{fig:06:tree}
\end{figure}

\subsection{窄通道规律}

把表 \ref{tab:06:pedigree} 与图 \ref{fig:06:tree} 合起来看，
可以提炼出一条本书称为\keyw{窄通道规律}的观察（\conf{M} 级，
由 \conf{C} 级语料归纳）：

\begin{proposition}[窄通道规律（教学命题）]\label{prop:06:channel}
  奶蛙变种的存活空间由两条约束夹出：
  其一，\textbf{可辨识下限}——丢失“黄色主体、大笑姿态、
  眼睛”等特征过多时，(C3) 失败，变种脱离谱系；
  其二，\textbf{新奇性下限}——与既有变种完全同质时，
  选择压力（社群的转发行为）不再偏好该变种。
  存活的变种集中于“轮廓趋同、细节发散”的窄通道内。
\end{proposition}

这条命题是“教学命题”：它把第五章
“变异率在‘可辨识’与‘足够怪’之间存在窄通道”的推论
落到了谱系表上。它的可检验形态是：
对任意时间窗内的奶蛙二创样本，
“轮廓特征保有率”的方差应显著小于
“细节特征保有率”的方差。
第十一章给出了这一检验的编码方案。
同时必须写明竞争性解释：轮廓趋同也可能仅仅是
AIGC 工具模板化输出的副产品（第八章），
而非社群选择的结果——两种解释在“工具固定、
社群不变”与“社群固定、工具变化”的对照样本上
可以区分。

\begin{nlcase}[案例 6-1：《植奶大战僵龙》与模因杂交]
  B 站 UP 主“暖气杀手 AAA”的《植奶大战僵龙》系列
  将奶蛙与《植物大战僵尸》（PVZ）相结合，
  制作多款高播放量视频，是奶蛙爆火的“强力引擎”之一
  \conf{C}（第一版教材口径）。
  从符号学上看，这是一次典型的\keyw{模因杂交}：
  奶蛙一方输出形象与动作程式（内容层嫁接），
  PVZ 一方输出玩法与场景框架（形式层嫁接），
  双方各保留自己“最硬”的维度。
  杂交体的归并判定：(C1) 对奶蛙侧成立（形体保留），
  (C2) 对两侧均成立（音轨、结尾套路 + PVZ 玩法），
  (C3) 由社群辨识实践确认（“植奶”“僵奶”成为独立词条）。
  杂交成功的前提是两个母体的形式层\keyw{不冲突}——
  这解释了为什么是塔防游戏而非叙事长篇
  更容易与短视频迷因杂交。
\end{nlcase}

\subsection{奶蛋：简化极限与“轮椅”叙事}

异形奶蛙中值得单独分析的是\term{奶蛋}：
奶蛙“简化到极致的产物”，在 B 站 UP 主“奶史\_”的
《选择你的奶伴》系列视频中高频出现，
与衍生语“安迪”结伴，因胜率极高而产生
“奶蛋轮椅”的说法（“安迪”空耳为“俺无敌”）
\conf{C}（第一版教材口径）。

奶蛋的学理价值在于它标定了谱系的\keyw{简化极限}：
当形体特征只剩“眼睛 + 大笑”，形象已经退化为一枚
近乎抽象的徽章——但只要它仍在既定形式
（音轨、结尾、衍生语）中出现，(C2) 就维持着簇归属。
“奶蛋轮椅”叙事则显示了内容与立场的耦合：
一个极简形象因为在特定叙事框架（对战、选择）
中“从无败绩”，被社群赋予了\keyw{强运}人设——
内容层的极度简化反而放大了立场层的投射空间。
这与表情包传播的一般规律一致：
形象越简单，可代入的自我越多。

\begin{nlcritique}[批判性反思：审丑与低质内容]
  奶蛙的“丑怪”不是中性事实。对这一谱系的批评集中在三点：
  其一，\textbf{审丑疲劳}——畸变与劣化作为刺激手段，
  在大量重复后会推高“出怪”的阈值，
  诱发更极端的形象变异，形成品质上的逐底竞争；
  其二，\textbf{低质内容的规模外部性}——《人民日报》报道
  明确把“用 AI 工具批量生成低质内容”列为官方反对的
  侵权形态\conf{A}\cite{jiang2026nailong}，
  而奶蛙谱系的生产方式恰好与“批量生成”高度重合
  （第八章详论）；
  其三，\textbf{对他人的刷屏外部性}——
  对不使用该迷因的人群，高密度的奶蛙内容
  构成信息流污染（第十章讨论校园场景中的这一问题）。
  本书把谱系当作符号系统分析，
  不等于为其中每一种用法背书。
\end{nlcritique}

\section{内容维度的正典侧：素材库与“糖”}\label{sec:06:canon}

奶蛙的内容维度不限于奶蛙自身：
正典奶龙的官方内容为二创提供了\keyw{素材库}。
这一节处理正典一侧的接受史（均为 \conf{C} 级，
据第一版教材记录）。

最受二创青睐的正典桥段是“我是奶龙，我才是奶龙”
（真假奶龙对峙）。据第一版教材的观察，
这一桥段受欢迎“并不是因为这一段好看，
而是因为这一段很糖”——
“糖”（社群用语）指内容幼稚、尬感强的特质
\conf{C}。社群甚至生产了讽刺性彩虹屁
（“我觉得真假奶龙这一段简直是智斗天花板……
我早已浑身湿透了”）来完成对“糖”的标记：
先模仿认真的夸奖语体，再让夸奖的内容
夸张到不可能被当真，从而实现反向评级。

围绕这一桥段形成了两条加工线：
其一，\keyw{倒放大法}——把音轨倒放，
“我是奶龙，我才是奶龙”变成
“我才是奶龙，欧蓝的手”之类的空耳文本
\conf{C}；其二，与特摄角色的\keyw{跨界绑定}
（如与“贝利亚”形象的并置），
使正典形象进入抽象语境。

\begin{nltermbox}{术语：空耳（soramimi）}
  源自日语“そらみみ”（空耳，幻听），
  指把一种语言的语音听成另一种语言
  （或本语言内其他词句）的再诠释实践。
  本书用法：\textbf{语音外壳保留、语义外壳被重新指派}
  的双层结构。在奶蛙语料中，
  空耳同时作用于正典素材（倒放产生的“欧蓝的手”）
  与衍生语（“安迪”$\leftrightarrow$“俺滴 / 俺无敌”）。
  空耳词等价项的判定标准与语音距离度量
  是第八章“可复现方法”的主题；
  此处只确立其符号学地位：空耳是\keyw{能指的劫持}，
  它让任何音轨都成为潜在的双关素材。
\end{nltermbox}

“糖”文化的存在解释了一个容易困惑的现象：
正典奶龙定位“治愈、欢乐”\conf{A}\cite{jiang2026nailong}，
其民间接受却是反讽式的——受众以“糖度”
为尺度重新分级官方内容，把最“幼儿向”的桥段
顶成传播峰值。接受（reception）与文本（text）
在这里显著脱钩：\keyw{民间评价体系
把正典当原料，而不当作品}。
这正是第二章“正典作为意义参照系”的具体运作方式——
社群以“糖”标记正典中“不可严肃对待”的部分，
为后续的畸变、拼贴清空了语义场地。
没有这场“清空”，奶蛙的畸变身体只会被视为
画坏了的奶龙；有了它，畸变才成为合法的风格。

\section{形式维度：叙事套路、音轨与衍生语}\label{sec:06:form}

\subsection{三件套：结构、音轨、语言}

奶蛙二创的形式层由三个高度稳定的构件组成。

第一，\keyw{叙事结构}：“压力桥段$\to$大笑收尾”。
典型视频以一个尴尬、失败、荒诞或高压的情境开场，
在情绪顶点处切入奶蛙的捧腹大笑，
以集体大笑为整个段落收尾\conf{C}\cite{shoutanji2026naiwa}。
这个结构让笑声承担“句号”的功能：
无论前文是什么，大笑都把它封存为“已经过去的事”。
第二，\keyw{音轨}：变调处理的爽朗笑声
（音源考据见第四章 \ref{sec:04:audio}），
二手评述还记录了其与变调版“不知火舞语音”的搭配
\conf{C}\cite{shoutanji2026naiwa}。
第三，\keyw{语言}：衍生语系统，
即下表的六个高频语音梗（均 \conf{C} 级，
据第一版教材整理）。

\begin{table}[htbp]\centering
  \caption{奶蛙衍生语系统（据第一版教材整理，\conf{C} 级）。
          “等效空耳”列即该词的再诠释空间，
          其判定标准第八章给出}
  \label{tab:06:langs}
  \small
  \begin{tabularx}{\textwidth}{@{}llX@{}}
    \toprule
    衍生语 & 主要运用场景 & 等效空耳（列举） \\ \midrule
    安迪 & 爬行奶蛙说话、奶蛋说话、奶蛙出场 &
      俺滴、安定、俺无敌、暗地 \\
    嘎嘎滴辣虾 & 运用场景广泛，高频说话效果 &
      刚刚点了谁、刚刚跌落水、刚刚掉了什 \\
    卡脚先 & 小剧场的转折或收尾 &
      刚救下、赶紧写、干啥先 \\
    无了 & 吃瘪、惊讶音效 & 哦了、好了 \\
    甘巴爹 & 肯定、疑问音效 & 干扒地、刚拔的 \\
    有用不 & 丰富语言，一般作“安迪”的下位替代 &
      忧郁菇、犹豫不 \\ \bottomrule
  \end{tabularx}
\end{table}

衍生语的使用有明确的\keyw{场景语法}：
“安迪”属于出场与移动（爬行到一半看向屏幕说“安迪”
是固定动作—语言组合），“无了”属于失败与惊愕，
“卡脚先”属于转折收束。用错场景不会被禁止，
但会被识别为外行——这正是立场维度
“圈内默示知识”的形式层证据。

\subsection{衍生动作：身体程式}

第一版教材记录的衍生动作有四种：\term{思考}、
\term{蹲起大笑}、\term{爬行}、\term{摆架势}
\conf{C}。其中爬行与“安迪”的绑定
（爬到一半看向屏幕、说出衍生语、继续爬行）
是动作—语言复合程式的典型；
摆架势则是“模仿奶蛙的青年们最常摆的动作”，
即程式从屏幕内迁移到屏幕外、
从形象动作迁移为使用者身体姿态的通道
（第十章的立牌合影、摆架势打卡是其空间化延伸）。

蹲起大笑另有一个“双人组合技”，
第一版教材以“画面过于神圣和有伤风化，
我们按下不表”为由未作展示\conf{C}。
本书沿用同一处理，并把它记为一个方法论样本：
\keyw{社群自身的语料边界实践}——
第一版教材的作者们已经在一个“什么不宜收录”的
问题上做了编辑决策。研究者的收录边界
（第十二章 E1--E3）与社群的自我边界不必一致，
但后者必须被记录，因为它是社群规范存在的直接证据。

\subsection{形式为什么稳定}

\begin{proposition}[形式稳定的三重来源（教学命题）]\label{prop:06:formstable}
  奶蛙形式层的稳定性（相较内容层的高替换率）
  可归因于三个互相加强的结构：
  (i) \textbf{学习成本}：形式构件数量少、组合规则简单，
  新用户一次接触即可习得；
  (ii) \textbf{识别成本}：固定音轨与结尾套路
  在信息流中提供即时可辨的“这是奶蛙视频”信号，
  满足平台分发对识别速度的要求；
  (iii) \textbf{组合开放}：内容层完全开放，
  任何素材都能接入同一形式，
  因此创新压力全部由内容层吸收，
  形式层无需变异即可持续产出“新”作品。
\end{proposition}

这一命题（\conf{M} 级）与第五章的
“变异—选择—保留”飞轮对接：形式层是被“保留”喂养的
惰性骨架，内容层是“变异”的活性表面。
第八章将从生产端给出同一结构的经济学版本：
AIGC 工具把内容变异的成本降到接近零，
而形式构件作为模板被工具固化，
于是“内容极活泼、形式极保守”成为均衡态。

\section{立场维度：用梗代替自述}\label{sec:06:stance}

\subsection{从表达到表态}

立场维度回答“用奶蛙在做什么”。本书归纳出
三层可观察的语用行为（\conf{M} 级归纳，
素材 \conf{C} 级）：

\begin{enumerate}[label=(\arabic*),leftmargin=2.4em]
  \item \textbf{自陈的替代}：用一张奶蛙梗图代替
        一句第一人称陈述，把“我很难受 / 我搞砸了 / 我绷不住了”
        翻译成“我们都这样”。戈夫曼的拟剧论
        把自我呈现分为前台与后台\cite{goffman1959presentation}；
        梗图在此充当\keyw{印象缓冲垫}——
        情绪信息被发出，但发出者退到了图像后面，
        既完成了表达，又规避了直接暴露的成本。
  \item \textbf{去抑制的通道}：网络环境的匿名性、
        不可见性与同步性缺失降低了表达门槛
        \cite{suler2004disinhibition}；
        “笑着崩溃”之所以可行，正因为笑声包裹了
        本该沉重的内容，使崩溃可以以娱乐的语体发出。
  \item \textbf{身份的默示声明}：正确使用衍生语、
        识别“过期”变体、掌握谱系细节，
        都是只有圈内成员能完成的操作。
        借用 Nissenbaum 与 Shifman 的框架，
        迷因使用是\keyw{可被社群审计的文化资本}
        \cite{nissenbaum2015contested}——
        转发什么、用什么语、对什么保持沉默，
        都在向同好发出可读取的信号。
        第七章的“逆反认证”机制（\ref{sec:07:reverse}）
        是这一逻辑在官方压制事件下的放大。
\end{enumerate}

\subsection{一份立场声明文本}

第一版教材的尾声收录了一篇社群信条式文本
（“奶蛙不叛”），其中罗列了奶蛙\keyw{不带来}的东西：
“没有勾心斗角”“没有拼死内卷，没有躺平摆烂”
“没有施加 PUA，没有拉升 BMI”“没有卡你 CET，
没有让你 EMO”\conf{C}（第一版教材原文节引）。

这份文本的符号学结构值得逐层拆解：
它的句式是排比否定（“没有 X”），
被否定的对象构成一份负面清单——内卷、绩点竞争、
身材焦虑、语言考试、情绪耗竭，
恰好是高校生活压力源的清单（第七章机制矩阵
“压力源”一列的民间版本）；
它把奶蛙定义为这些压力的\keyw{缺席者}，
并以此号召“发自内心地……笑出来吧”。
换言之，社群自己写下的立场声明
把“捧腹大笑”明确登记为\keyw{对压力情境的应答}，
而不是随机娱乐。这份文本是 \conf{C} 级材料，
不能当作“大学生使用奶蛙的动机”的证据；
但它至少证明：\keyw{在社群的自我理解中，
奶蛙与压力话语被组织进了同一个意义结构}。
动机是否果真如此，是第十一章问卷要回答的问题。

\begin{nlmethod}[研究方法提示：立场文本的取样]
  分析“社群立场”时，警惕以声明文本代替行为证据：
  信条式文本是\keyw{自我表述}，
  存在美化、跟风与戏仿三层混杂
  （“奶蛙不叛”本身是否也是玩梗？）。
  可行的取样纪律：
  (1) 声明文本只用于界定“社群自我理解”，
  不用于推断个体动机；
  (2) 行为证据（何时、对谁、在什么情境使用）
  优先于文本证据；
  (3) 对声明文本本身做语体标注——
  宗教戏仿、宣言戏仿的语体标记
  恰恰说明作者清楚这不是一份严肃文本。
\end{nlmethod}

\section{三维的不对称：“一软两硬”}\label{sec:06:asym}

本章的分析可以收拢为一个结构性判断：

\begin{table}[htbp]\centering
  \caption{奶蛙符号系统的“一软两硬”结构（本书归纳，\conf{M} 级）}
  \label{tab:06:asym}
  \small
  \begin{tabularx}{\textwidth}{@{}llX@{}}
    \toprule
    维度 & 状态 & 直接表现 \\ \midrule
    内容 & 软（高替换） &
      形象谱系横跨物种与器物；
      任何“黄色大笑生物”皆可接入；
      正典桥段可反讽征用 \\
    形式 & 硬（高稳定） &
      结尾集体大笑、变调音轨、
      衍生语场景语法、压力桥段结构长期不变 \\
    立场 & 硬（高约定） &
      “绷不住但用笑处理”的解读框架稳定；
      场景语法错误会被识别为外行；
      官方否认反而加固圈内立场 \\ \bottomrule
  \end{tabularx}
\end{table}

“一软两硬”解释了奶蛙作为符号系统的生存策略：
\keyw{用最硬的形式承载最软的内容}。
形式与立场的稳定保证了辨识与归属（(C2)(C3)），
内容的开放保证了变异供给（(C1) 的宽容度），
三者合起来正好是一个“可无限换壳的仪式框架”。
巴赫金笔下的狂欢节正是类似的结构：
仪式形式年年相同，插科打诨的内容随时更新
\cite{bakhtin1968rabelais}。
区别在于媒介条件——狂欢节一年一度，
奶蛙一天可以刷到几十次，
仪式框架的重复频率被平台分发放大了几个数量级
（第七章将讨论这种高频重复与情绪耗竭的关系）。

由此回到章首的承诺：为什么说奶蛙是“三个层面的叠加”
（第一章）而不仅是“一个梗”——
因为它的符号系统本身就分层运行：
正典素材库（内容原料）、
民间谱系（内容产品）、
形式—立场复合体（仪式框架）。
第九章讨论产业时将看到，
这个结构还决定了它的经济形态：
内容层无人收费，形式层无法收费，
最终被商业捕捉的是立场的周边化（立牌、玩偶）。

\section{本章小结}\label{sec:06:summary}

本章用 Shifman 三维框架解剖了奶蛙语料。
内容维度给出三级谱系（换皮、异形、客串）
与简化极限（奶蛋），并提出窄通道规律
（轮廓趋同、细节发散，\conf{M} 级）；
正典侧的“糖”文化与倒放空耳说明了素材库的
反讽式接受机制。形式维度拆解了
“结构—音轨—语言”三件套与四种衍生动作，
并给出形式稳定的三重来源命题。
立场维度区分了自陈替代、去抑制通道、
身份默示声明三种语用行为，
并以“奶蛙不叛”文本展示了社群的自我理解。
总的判断是“一软两硬”：
内容高替换、形式高稳定、立场高约定化。
遗留问题有二：
其一，形式层为什么在 AIGC 时代格外稳定，
需要生产端的证据（第八章）；
其二，立场层的三个语用行为
是否真的对应使用者的心理过程，
符号分析无法回答，需要受体研究（第十一章）。

\begin{nlexercises}
  \item 取表 \ref{tab:06:pedigree} 中任意三个变种，
        逐条检验归并标准 (C1)(C2)(C3)，
        并指出哪一条标准在判定中起了决定作用。
        若三个变种中有一个你认为应当\keyw{移出}谱系，
        写出你的理由（丢失了哪个不可再丢的特征）。
  \item “糖”是社群的负面—反讽评级词。
        再找一个你所知的网络社群评级词
        （如“雷”“尬”“典”），
        写出它的判定操作（什么可观察特征使一段内容获得该评级），
        并讨论：当评级词本身被官方或广告征用后，
        评级效力会发生什么变化？
  \item 本章把“奶蛙不叛”文本定位为
        “社群自我理解”而非“个体动机证据”。
        设计一个区分二者的最小研究方案：
        至少一个行为指标、一个文本指标，
        并说明两者结果不一致时你接受哪一个、为什么。
\end{nlexercises}

\nlkeywords 三维框架、形象谱系、换皮奶蛙、异形奶蛙、奶蛋、
空耳（soramimi）、糖、倒放大法、衍生语、衍生动作、
场景语法、一软两硬、窄通道规律
